\title{LongRoPE: Extending LLM Context Window Beyond 2 Million Tokens}

\begin{abstract}
Expanding the context window in large language models (LLMs) is highly sought after. Nevertheless, existing context window extensions typically cap at approximately 128k tokens, largely due to substantial fine-tuning expenses, limited availability of sufficiently lengthy textual data, and destabilizing effects from introducing novel token positions.

In this work, we present LongRoPE, a method that, for the first time, expands the context window of pre-trained LLMs to an extraordinary \textbf{2048k} tokens, requiring at most 1k fine-tuning steps on sequences up to 256k in length, while preserving performance at the standard short context range. Our approach relies on three primary innovations: (i) we uncover and leverage two specific types of non-uniformities in positional interpolation through an efficient search procedure, yielding improved initialization for fine-tuning and supporting an 8$\times$ extension without fine-tuning; (ii) we propose a progressive extension process, wherein we first fine-tune an LLM on 256k tokens before applying a second round of positional interpolation to the resulting extended model, thereby attaining a 2048k context window; and (iii) we retune LongRoPE at a context length of 8k to restore optimal short-range performance. Comprehensive experiments utilizing LLaMA2 and Mistral on diverse tasks confirm the superiority of our technique. Importantly, LongRoPE-enhanced models retain their original structure except for minimal adjustments to the positional embeddings, and can take advantage of most existing optimizations. Code for our method will be provided at \texttt{https://github.com/microsoft/LongRoPE}
\end{abstract}

\section{Introduction}

Although Large Language Models (LLMs) have achieved impressive results across a range of tasks~\cite{gpt4,llama2}, they are constrained by a fixed context window, such as the 4096-token restriction in LLaMA2~\cite{llama2}. When input sequences extend beyond this limit, model accuracy deteriorates because the LLM encounters position indices absent from its training regime. This constraint hinders the effective use of LLMs in critical applications, including in-context learning scenarios involving extensive example sets~\cite{Huang2023FewerIM} and the operation of LLM agents~\cite{park2023generative,madaan2023selfrefine}.

Recent studies have demonstrated that it is feasible to push the context window of a pre-trained LLM to approximately 128k tokens by applying fine-tuning techniques using longer documents~\cite{longlora,pi,yarn,zhang2024soaring,liu2023scaling}. Despite this progress, three key challenges remain in the pursuit of even larger context windows. Firstly, assigning uninitialized position indices for new token positions results in a large influx of extreme values, which causes distribution shifts and severely hampers fine-tuning convergence~\cite{pi}. This problem is intensified in scenarios where the jump is substantial—such as an expansion from 4k to over 1000k tokens, which involves introducing more than 90\% entirely novel position indices. Secondly, it is generally necessary to provide training samples that match the target sequence length for effective fine-tuning. Yet, contemporary corpora contain very few instances of texts surpassing 1000k tokens. Furthermore, the computational burden of processing such extra-long sequences is immense, demanding excessive training durations and large-scale GPU allocations. Thirdly, scaling up the context window to extreme sizes causes the model’s attention to be allocated across a vast number of positions, resulting in diluted attention weights and reduced performance when dealing with shorter, original-context tasks~\cite{pi}.

A commonly used strategy to address the initial challenge involves adapting RoPE positional embeddings~\cite{rope,pi} by mapping new position indices back to the original pretrained range, as illustrated in Fig.~\ref{fig:overview}. Specifically, Position Interpolation (PI)~\cite{pi} modifies RoPE’s rotary angles using a linear scaling based on the extension ratio. In contrast, NTK~\cite{ntk,dynamicntk} proposes non-uniform methods of interpolation and extrapolation tailored to different dimensions of RoPE. Meanwhile, YaRN~\cite{yarn} partitions RoPE dimensions into three distinct groups according to frequency, applying extrapolation, NTK interpolation, and linear interpolation methods to each category in turn. Despite these advancements, the positional embeddings within Transformer models inherently display \emph{complex non-uniform information entropy}. Current methods do not sufficiently exploit this nuanced irregularity, which leads to the potential loss of information and ultimately curtails the possible expansion of the context window.

Section~\ref{sec:analysis} empirically demonstrates two main insights: \textbf{(1)} Effective positional interpolation must address two sources of non-uniformity—differences across RoPE dimensions and variability in token positions. Lower RoPE dimensions and tokens appearing earlier in the sequence require less interpolation. However, the optimal configuration is influenced by the specific target extension length. \textbf{(2)} Incorporating awareness of these non-uniformities into the interpolation process helps to better preserve the representational fidelity of the original RoPE, particularly for crucial dimensions and token positions. This approach reduces information loss resulting from interpolation, thereby yielding superior initialization for subsequent fine-tuning. Furthermore, it facilitates extending the sequence length by up to 8$\times$ even in the absence of further fine-tuning.

Building upon these results, we propose \textbf{LongRoPE}, a robust technique designed to extend the context window of LLMs to over 2 \emph{million} tokens. LongRoPE incorporates three principal innovations. To begin, {LongRoPE} leverages complex multidimensional variations in positional interpolation, discovering optimal rescaling coefficients for each RoPE dimension’s rotation angles as a function of token positions. Given that the parameter search space for these rescale factors grows exponentially relative to the desired context length expansion, {LongRoPE} adopts an evolutionary search algorithm enhanced with two optimization strategies, thereby significantly increasing search efficiency. An illustration of the resulting rescaled RoPE can be found in Fig.~\ref{fig:overview}.

Subsequently, {LongRoPE} employs a resource-efficient and iterative extension approach to realize a 2048k context window, circumventing the necessity for explicit fine-tuning on extremely lengthy text corpora, which are rare and often unavailable. This process initiates by identifying a viable 256k context length for the pre-trained LLM and applying fine-tuning at this intermediate length. Leveraging our non-uniform positional interpolation, which facilitates an 8$\times$ context window expansion even without additional fine-tuning, a second phase is undertaken wherein updated RoPE rescale factors are searched for using the previously fine-tuned extended LLM. This comprehensive approach results in the successful extension of the context window to 2048k for both LLaMA2 and Mistral~\cite{mistral}.

To address the decline in performance when handling the original, shorter context window, {LongRoPE} applies further modifications to the RoPE rescale factors within the expanded LLM. Following the approach taken when scaling from 256k to 2048k, we now reduce the context windows to 4k and 8k on the LLM fine-tuned with a 256k context. Our search algorithm is leveraged in this process to minimize positional interpolation effects. For inference, if the input sequence is shorter than 8k, RoPE is recalibrated using the rescale factors determined through this search.

A comprehensive series of experiments involving multiple LLMs and diverse long-context benchmarks validates the efficacy of our approach. Our results indicate that LongRoPE consistently sustains low perplexity scores over evaluation lengths ranging from 4k to 2048k, attains passkey retrieval accuracies exceeding 90\%, and maintains accuracy on standard tasks formulated within a 4096-token context window that is on par with existing methods. Furthermore, LongRoPE is compatible with all LLMs utilizing RoPE embeddings. We plan to make both our implementation and the LongRoPE-2048k models publicly available.

\section{Non-uniformity in Positional Interpolation}
\label{sec:analysis}

\subsection{Preliminary}

Transformer models require explicit positional information, often in the form of position embedding, to represent the order of input tokens. Our work focuses on the RoPE~\cite{rope} position embedding, which is widely used in recent  LLMs. For a token at position index $n$, its corresponding RoPE encoding can be simplified as follows:

\begin{equation}
		\fontsize{7.8}{7.8} \selectfont
	\label{eq:rope}
	\left[ cos(n\theta_0), sin(n\theta_0), cos(n\theta_1), \cdot\cdot\cdot, cos(n\theta_{d/2-1}), sin(n\theta_{d/2-1}) \right]   
\end{equation}

where $d$ denotes the dimensionality of the embeddings, $n\theta_i$ indicates the rotary phase for a token at position $n$, and $\theta_i=\theta^{-2i/d}$ specifies the rotation frequency parameters. In RoPE, $\theta$ is typically set to a base value of 10000.

\textbf{\emph{Context window extension ratio $s$} and positional interpolation}. The parameter $s$ is introduced as the ratio between the new, extended context length $L'$ and the original context length $L$, that is, $s=\frac{L'}{L}$.

To extend the context window from $L$ to $L'$, current positional interpolation methods suggest downscaling rotation frequencies $\theta_i$ based on extension ratio $s$. Let $\beta=\theta^{2/d}$, and $\lambda$ denote the actual rescale factor related to $s$, we   unify these positional interpolation methods as follows:

\begin{equation}	
	\fontsize{7.5}{7.5} \selectfont
	\label{eq:unified_rope}
	\begin{aligned}
		\left[cos\left(\frac{n}{\lambda(\beta)^0}\right), sin \left(\frac{n}{\lambda(\beta)^0}\right), cos\left(\frac{n}{\lambda(\beta)^1}\right), \cdot\cdot\cdot,  sin \left(\frac{n}{\lambda(\beta)^{d/2-1}}\right) \right]   
	\end{aligned}
\end{equation}

\textbf{Linear positional interpolation (PI)}. PI~\cite{pi} proposes using linear interpolation on position indices when operating within the pre-trained sequence length. When extending to a new target length with ratio $s$, the rotational angles assigned to each position are uniformly scaled down by a factor of $\lambda=s$ across the entirety of the RoPE dimensions. Despite this straightforward adaptation, this approach compresses the positional information, which results in diminished capacity for the model to discern between tokens that are near each other. Consequently, the effectiveness of PI decreases notably as the extension ratio becomes larger.

\textbf{NTK-based interpolation and extrapolation}. ~\cite{ntk,dynamicntk} approach RoPE through the lens of information encoding, leveraging Neural Tangent Kernel (NTK) theory~\cite{jacot2018neural,tancik2020fourier}. To address the problem of overcrowded positions encountered in PI, these works propose a strategy of redistributing the interpolation demand among RoPE’s dimensions. This method reduces the scaling of lower (high frequency) dimensions while increasing it for higher (low frequency) dimensions, thereby enabling both positional interpolation and extrapolation with $\lambda=s^{i}$. The refined dynamic NTK approach~\cite{dynamicntk} further adapts the extension factor at each position in accordance with the sequence length. Distinct from PI, which requires fine-tuning, NTK-oriented techniques are capable of extending the context window even without additional fine-tuning, although practical extension is often capped at 4$\times$ the original window.

\textbf{YaRN}~\cite{yarn} presents notable advancements in the effectiveness of positional interpolation. The method segments the RoPE dimensions into three categories according to their frequency characteristics, applying distinct interpolation techniques to each. Dimensions with high frequencies are treated with extrapolation ($\lambda$=1), whereas those with low frequencies utilize linear interpolation (PI). Intermediate-frequency RoPE dimensions adopt the NTK approach. A central aspect of YaRN is its partitioning strategy for RoPE dimensions; however, this strategy is presently determined by manual empirical evaluation, which may lead to less-than-optimal results when applied to newly developed LLMs.

\subsection{Study on Non-uniform Positional Interpolation}

Building on insights from NTK and YaRN, we observe that their improvements stem from the introduction of non-linearity, particularly by leveraging varying frequencies in the RoPE dimensions to tailor interpolation and extrapolation. Despite these advancements, the prevailing non-linear methods are still predominantly based on heuristics crafted by humans. This leads us to consider two central questions: (1) Does the present approach to positional interpolation represent the best possible method? (2) Are there additional non-linear forms that have yet to be discovered or utilized?

In addressing these questions, we employ evolution search (refer to Sec.~\ref{sec:method}) to identify improved non-uniform positional interpolations for LLaMA2-7B. The evolution process is steered by evaluating perplexity, leveraging five randomly chosen samples from the PG19~\cite{pg19} validation set. Our experiments yield several notable insights, detailed below.

\textbf{Finding 1}: \textit{RoPE dimensions exhibit substantial non-uniformities, which are not effectively handled by current positional interpolation methods.}

We optimize $\lambda$ for each RoPE dimension as specified in Eq.~\ref{eq:unified_rope}. Table~\ref{tbl:exp1} presents the perplexity results of LLaMA2-7B across different approaches on the PG19 and Proof-pile~\cite{proof-pile} test sets, all evaluated without any fine-tuning. The solution identified through our search process yields substantial improvements, indicating that prevailing linear (PI) and non-uniform (Dynamic-NTK and YaRN) interpolation methods do not achieve optimal performance. It is particularly noteworthy that YaRN performs worse than both PI and NTK on PG19; for non-fine-tuned large language models, YaRN fails to attain the designated context window size. Specifically, with an 8k context window, YaRN's perplexity rises sharply after processing 7k tokens.

Our investigation reveals that the rescaling coefficients $\lambda$ in Eq.~\ref{eq:unified_rope} are distributed unevenly, in contrast to the constant scale $s$ utilized in PI, the NTK formulation, or the group-based approach in YaRN. Employing these heterogeneous scaling factors leads to notable enhancements in LLaMA2’s language modeling (as measured by perplexity) across 8k and 16k token context lengths, even in the absence of additional fine-tuning. This improvement can be attributed to the resultant positional encodings, which more faithfully retain the characteristics of the original RoPE—particularly for important dimensions—thereby alleviating the model’s challenge in discriminating among tokens with proximal positions.

\textbf{Finding 2}: \textit{RoPE for the initial tokens in the input sequence should be extrapolated with less interpolation.} 

We propose that for the first $\hat{n}$ tokens of input sequences, the degree of interpolation performed by their RoPE embeddings should be minimized. This suggestion is motivated by evidence that these tokens, receiving comparatively high attention scores, play a pivotal role in attention mechanisms—an effect also identified by Streaming LLM~\cite{streamingllm} and LM-Infinite~\cite{han2023lminfinite}. To examine this hypothesis, we enlarge the context window to 8k and 16k with both PI and NTK methods, but exclude the initial $\hat n$ (0,2, ..., 256) tokens from any interpolation. Notably, if $\hat n$=0, this corresponds to standard PI and NTK implementations. Our results, summarized in Table~\ref{tbl:tokenposition}, yield two major insights: \textbf{(1)} Performance increases when interpolation is omitted for the leading tokens in a sequence. \textbf{(2)} The best value for $\hat n$—the count of such initial tokens—varies with the desired extension length.

\textbf{Finding 3}: \textit{Non-uniform positional interpolation effectively extends LLM context window in both fine-tuning and non-fine-tuning settings.} 

Our results indicate that employing our optimized non-uniform position interpolation substantially enhances extension capabilities at 8k and 16k token lengths even in the absence of fine-tuning. However, scaling to larger context sizes necessitates fine-tuning. Consequently, we further fine-tune LLaMA2-7B leveraging our optimized RoPE parameters for a 64k-token context window (see Appendix for detailed configurations). As evidenced in Table~\ref{tbl:finetune}, this strategy enables our approach to consistently outperform both PI and YaRN—across both pre- and post-fine-tuning scenarios. The improvement is attributed to the tailored use of non-uniform positional interpolation, which limits information degradation and yields a superior initialization prior to fine-tuning.

\textbf{Summary}. Our investigation identifies two sources of non-uniformity: heterogeneous RoPE dimension scaling and uneven distribution of token positions. Harnessing these aspects within positional interpolation drives significant gains in the ability of LLMs to handle extended contexts.

\section{LongRoPE}

\label{sec:method}
Building on these insights, we propose LongRoPE, which initially incorporates an optimized search strategy designed to leverage both types of non-uniformity. Utilizing this approach, LongRoPE successfully expands the context window for LLMs to more than 2 million tokens.

\subsection{Problem Formulation}

These two sources of non-uniformity result in an extensive solution space and complicate the optimization process. To tackle this challenge, we reformulate the multidimensional non-uniform position interpolation optimization as a search problem.

For a LLM targeting a  context window size of $L'$ and lengthy input documents $\mathbf{X}$, where each $\mathbf{x}\in\mathbf{X}$ surpasses $L'$ in token length, we denote the original rotary angle of the $i^{th}$ dimension in RoPE embedding at token position $n$ as  $\frac{n}{\beta^i}$. The optimization problem is then formulated as follows:

\begin{equation}
\begin{gathered}
		\label{eq:problem}

\underset {\mathbf{x}\in\mathbf{X}; \, |\mathbf{x}|\ge L'}{ \operatorname {arg\,min} } \,  \mathcal{L}\left( \text{LLM}\left(\text{RoPE}, \mathbf{X}\right)\right), \;\text{where}
\\
\scalebox{0.93}{
$\underset { \begin{subarray}{l} i=0,\ldots,\frac{d}{2}-1; \\ n\in[0,|\mathbf{x}|); \end{subarray} }{ \text{RoPE}_{(n)} }
=
{\left[\ldots, \cos\left(\mathbb{I}(\hat{\lambda}_i, \hat{n}) \frac{n}{\beta^i} \right), \; \sin\left(\mathbb{I}(\hat{\lambda}_i, \hat{n}) \frac{n}{\beta^i} \right), \ldots\right]}$
}
\\
\text{where } \mathbb{I}(\hat{\lambda}_i, \hat{n}) = \begin{cases}
1 & \text{if } n < \hat{n} \\
{\frac{1}{\lambda}_i} & \text{if } n \geq \hat{n}
\end{cases}
\end{gathered}
\end{equation}

In this framework, we introduce rescale factors $\mathbb{I}(\hat{\lambda}_{i},\hat{n})$ to address both types of non-uniformity. Here, $\hat{\lambda_i}$ represents dimension-based non-uniformity in RoPE, while $\hat{n}$ captures the non-uniformity relative to token positions. The function $\mathbb{I}(\hat{\lambda}_{i},\hat{n})$ adjusts the rotational angle associated with the $i^{th}$ dimension of RoPE: $\hat\lambda_{i}$ serves as the scaling coefficient, and $\hat{n}$ defines the positional threshold. For token positions less than $\hat{n}$, the original rotary positional embedding $\frac{n}{\beta^i}$ remains unaltered, as the rescaling is inactive. Once token positions reach $n \ge \hat{n}$, the rescaling factor is applied accordingly.

Assuming a target context window length of $L'$, our goal is to determine the most suitable rescale coefficients ($\mathbb{I}(\hat{\lambda}_{0},\hat{n})$, $\mathbb{I}(\hat{\lambda}_{1},\hat{n})$, ... ,$\mathbb{I}(\hat{\lambda}_{i},\hat{n})$, ...) corresponding to each RoPE dimension from the $1^{st}$ to the $d^{th}$. With these adjusted $\text{RoPE}$ parameters, the target $\text{LLM}$ is expected to minimize the next token prediction loss, denoted by $\mathcal{L}$ (which correlates with perplexity), when provided with input sequences $\mathbf{X}$ whose token counts equal $L'$.

\subsection{Searching the Non-uniform Position Interpolation}

In order to address the challenge outlined in Eq.~\ref{eq:problem}, we propose a straightforward but powerful approach that identifies the best RoPE rescale factors, enabling comprehensive utilization of the complex, multidimensional variations present in position embedding.

\textbf{Search space}. We construct an extensive search space to account for both types of non-uniformity. The search space is outlined in Table~\ref{tbl:searchspace}. In particular, our method permits individualized rescale factors for each RoPE embedding dimension. To streamline the search space configuration, we optimize over $\lambda_i$ and $\hat{n}$ rather than directly optimizing over $\mathbb{I}(\hat{\lambda}_{i},\hat{n})$, with $\hat{\lambda}_{i}=1/\lambda_i$. According to Table~\ref{tbl:searchspace}, $\lambda_i$ can be selected within the range starting at 1.0 (representing straightforward extrapolation) and extending up to $s\times1.25$ (representing an interpolation range surpassing PI), incremented by steps of 0.01, where $s$ denotes the intended expansion ratio of the context window.

The parameter $\hat{n}$ specifies how many of the earliest token positions preserve their original configuration, foregoing position interpolation and maintaining the initial RoPE embedding. In our experiments, we permit $\hat{n}$ to take values from the set \{0, 1, 2, 4, 8, 12, 16, 20, 24, 28, 32, 64, 128, 256\}. Setting $\hat{n}=0$ means that every token position adopts the optimized rescale factors.

\textbf{Evolution-based search}. The search space described in Table~\ref{tbl:searchspace} comprises a vast array of positional interpolation options, making efficient search highly nontrivial. For instance, when the extension ratio is $s=4\times$, the total number of combinations reaches $400^{128/2}\times14 = 4\times10^{167}$, and this figure grows exponentially as the extension ratio increases. To efficiently navigate this enormous space, we adopt evolution search~\cite{guo2020single} and propose two optimization strategies that substantially enhance the search process. The overall procedure is outlined in Algorithm~\ref{alg:search}.

\textit{Optimized initial population generation}. In place of creating all $P$ rescale factors by random assignment, our approach includes the three RoPE rescale factors for PI, NTK, and YaRN directly as members of the initial population. The rest of the individuals, numbering $P-3$, are produced by subjecting these three rescale factors to random mutations, each applied with a probability of $p$.

\textit{Monotonically non-decreasing constraint}. Following the creation of the initial population, we evaluate the perplexity of each candidate using the target LLM. This involves incorporating the respective RoPE rescale factors into the model and then measuring perplexity on the input $\mathbf{X}$. The most promising $k$ individuals—those with the lowest perplexity—are selected as parents for the next evolutionary step. Nevertheless, the immense search space means that straightforward applications of mutation and crossover may frequently result in suboptimal candidates, prompting redundant perplexity assessments. This challenge becomes more pronounced when $L'$ is high, as each perplexity inference demands considerable computation time.

To tackle this issue, we enforce a monotonic non-decreasing constraint on the sequence of sampled RoPE rescaling factors, such that $\lambda_i \le \lambda_{i+1}$ for all $i$. Only RoPE parameter sets that adhere to this requirement are utilized in LLM perplexity assessments, which leads to a considerable reduction in search overhead. More concretely, this means that the rescaling factors $\lambda_i$ must be monotonically increasing with respect to the RoPE dimension (for $i = 0,\ldots,63$). This approach draws inspiration from NTK theory~\cite{jacot2018neural,tancik2020fourier,ntk}, which indicates that lower-indexed (higher-frequency) dimensions require less interpolation (manifested as smaller values of $\lambda_i$), while higher-indexed (lower-frequency) dimensions are capable of greater interpolation (corresponding to larger $\lambda_i$ values).

\begin{algorithm}[t]
	\small
	\caption{The search algorithm}
	\textbf{Input:} target LLM, input samples $\mathbf{X}$, population size $P$, mutation size $N_1$, crossover size $N_2$, maximum iterations $\mathcal{T}$, mutation probability $p$\\

	\begin{algorithmic}[1]
		\label{alg:search}
		\STATE Initialize Top-k as $\phi$;	
		\STATE $\text{P}_0$ = \textit{Initialize\_population\_with\_optimization} ($P$, $\mathbf{X}$, $p$); 
		\FOR{$i$ = $1$ to $\mathcal{T}$}
		\STATE \textit{Compute\_perplexity} (LLM, $\text{P}_{i-1}$, $\mathbf{X}$);
		\STATE Top-k = \textit{Update\_Topk} (Top-k, $\text{P}_{i-1}$);
		\STATE $\text{P}_{mutation}$ = \textit{Mutation\_with\_mono\_constraint} (Top-k, $N_1$, $p$); 
		\STATE $\text{P}_{crossover}$ = \textit{Crossover\_with\_mono\_constraint} (Top-k, $N_2$);
		\STATE $\text{P}_i$ = $\text{P}_{mutation}$ $\cup$ $\text{P}_{crossover}$ $\cup$ Top-k;
		\ENDFOR
		\STATE Output the element in Top-k exhibiting the minimum perplexity;
	\end{algorithmic}
\end{algorithm}

\textbf{8$\times$ extension without fine-tuning}. Through evolutionary search, our approach efficiently determines optimal non-uniform RoPE rescale factors, retaining crucial dimensions and positions to reduce information loss due to interpolation. As shown in Fig.\ref{fig:non-ft}, this strategy enables the LLaMA2 model to expand its context window from 4k to 32k tokens without the need for fine-tuning. By comparison, current techniques including PI as well as non-uniform NTK and YaRN result in substantial increases in perplexity once the context is extended beyond 2$\times$.

\subsection{Extending LLM Context Window to 2048K}
\label{sec:extendto2048k}

\textbf{Progressive extension to 2048k}. In this section, we present our approach for expanding the context length of pre-trained LLMs from the conventional 4k tokens to more than 2048k tokens. Our non-uniform positional interpolation is capable of achieving an extension up to 8$\times$ the original length without the need for additional fine-tuning. However, for far larger expansions, such as extending by a factor of 512$\times$, fine-tuning becomes essential. A potential strategy involves identifying suitable RoPE rescale parameters specific to the desired 2048k window before initiating fine-tuning. Nonetheless, this process is hampered by the immense computational costs involved. Furthermore, our empirical investigations suggest that fine-tuning LLMs at such large extension ratios is inherently difficult (refer to Appendix for further discussion).

Fortunately, LongRoPE demonstrates strong performance for both base and fine-tuned variants of extended LLMs. In this context, we propose an efficient and incremental approach that attains the target of 2048k using merely 1k fine-tuning iterations, all within a training sequence length capped at 256k.

$\Diamond$ \textit{Applying LongRoPE search to extend pre-trained LLMs to 256k}. Using LLaMA2 as a case study, we explore scaling to context windows of 128k and 256k tokens. This process results in expansion factors of 32$\times$ and 64$\times$, respectively.

$\Diamond$ \textit{Fine-tuning to 256k}. Subsequently, we adapt the pre-trained LLM for a context window of 256k. To accomplish this, we initially fine-tune LLaMA2 over 400 steps utilizing RoPE rescaling parameters set for 128k. Following this phase, we substitute the RoPE rescaled factors with those configured for 256k in the resulting checkpoint, and proceed with an extra 600 fine-tuning steps. This staged approach demonstrates greater efficiency compared to performing all fine-tuning directly at 256k.

$\Diamond$ \textit{Scaling fine-tuned extended LLM to 2048k using LongRoPE search}. In the final step, we conduct an additional search phase on the LLM previously fine-tuned for a 256k context length. Through this process, we achieve a substantial context window size of 2048k, all without the need for extra fine-tuning. Consequently, the model attains a total extension factor of $512\times$.

\textbf{Mitigation of performance degradation in shorter context windows}.  When scaling up the context window to a substantial 2048k length, we observe that model effectiveness deteriorates in the initial, shorter context span. This phenomenon has been reported previously for positional interpolation~\cite{pi}: constraining position embeddings to fit a greatly enlarged context range compresses the embeddings in lower-indexed positions, which in turn hampers model performance. For example, with an extension ratio of 512$\times$, the available embedding space for the initial 4k positions is heavily compressed, resulting in substantial overlap and crowding.

To address this issue, we conduct an additional evolution search on the expanded LLM, optimizing the RoPE rescale factors for shorter context lengths, such as 4k and 8k. The upper limit for the searched $\lambda$ is lowered, as shorter contexts require reduced positional interpolation. At inference time, the LLM automatically updates the respective RoPE rescale factors according to the input.

\section{LongRoPE}

\label{sec:method}
Building on these observations, we propose LongRoPE, a method that initially employs an efficient search algorithm designed to leverage the two distinct non-uniformities. Utilizing this approach, LongRoPE is able to scale the context window of LLMs to exceed 2 million tokens.

\subsection{Problem Formulation}

These two forms of non-uniformity expand the solution space considerably and add layers of complexity to the optimization process. To tackle this issue, we recast the multidimensional non-uniform position interpolation optimization challenge into a search problem.

For a LLM targeting a  context window size of $L'$ and lengthy input documents $\mathbf{X}$, where each $\mathbf{x}\in\mathbf{X}$ surpasses $L'$ in token length, we denote the original rotary angle of the $i^{th}$ dimension in RoPE embedding at token position $n$ as  $\frac{n}{\beta^i}$. The optimization problem is then formulated as follows:

\begin{equation}
\begin{gathered}
		\label{eq:problem}

\underset {\mathbf{x}\in\mathbf{X}; \, |\mathbf{x}|\ge L'}{ \operatorname {arg\,min} } \,  \mathcal{L}\, \left( \text{LLM} (\text{RoPE}, \mathbf{X})\right),	\text{with } 
\\
\scalebox{0.93}{
$\underset {\begin{subarray}{l} i=0,\cdot\cdot,\frac{d}{2}-1;\\ n\in[ 0,|\mathbf{x}|);\end{subarray}}{\text{RoPE}_(n)}= {\left[\cdot\cdot,\cos\left(\mathbb{I}(\hat{\lambda}_{i}, \hat{n})\cdot\frac{n}{\beta^i}\right), \sin\left(\mathbb{I}(\hat{\lambda}_{i}, \hat{n})\cdot\frac{n}{\beta^i}\right),\cdot\cdot\right] }$}\\
	and \mathbb{I}(\hat{\lambda}_{i},\hat{n})=\begin{cases}
	1&\mbox{\; if $n<\hat{n}$}\\
{\frac{1}{\lambda}_{i}}&\mbox{\; if $n\geq\hat{n}$}
\end{cases}
	\end{gathered}
\end{equation}

In this approach, we define a collection of rescale factors, $\mathbb{I}(\hat{\lambda}_{i},\hat{n})$, intended to address two distinct types of non-uniformities. Here, $\hat{\lambda}_i$ represents discrepancies across RoPE dimensions, while $\hat{n}$ captures irregularities associated with token positions. Concretely, for the $i^{th}$ RoPE dimension, $\mathbb{I}(\hat{\lambda}_{i},\hat{n})$ serves to adjust the corresponding rotation angle; $\hat{\lambda}_{i}$ acts as the scaling coefficient and $\hat{n}$ specifies the token position threshold. Up to the first $\hat{n}$-1 tokens, this adjustment is not applied, so the unmodified RoPE rotary angle $\frac{n}{\beta^i}$ remains in effect. Once the position $n$ meets or exceeds $\hat{n}$, the scaling factor is incorporated into the computation.

Assuming a desired context window of size $L'$, the aim is to determine the best rescale factors ($\mathbb{I}(\hat{\lambda}_{0},\hat{n})$, $\mathbb{I}(\hat{\lambda}_{1},\hat{n})$, ...$\mathbb{I}(\hat{\lambda}_{i},\hat{n})$...) for each RoPE dimension from the $1^{st}$ up to the $d^{th}$. With these optimized scaling values applied to $\text{RoPE}$ in the target $\text{LLM}$, the system should minimize the next-token prediction loss $\mathcal{L}$ (i.e., reduce perplexity) on input sequences $\mathbf{X}$ of length $L'$.

\subsection{Searching the Non-uniform Position Interpolation}

To address the issue outlined in Eq.~\ref{eq:problem}, we propose a straightforward but powerful approach that identifies the best RoPE rescale factors, thereby enabling the model to make full use of the complex, multidimensional variations present in position embedding.

\textbf{Search space}. We construct an extensive search space that accommodates the two forms of non-uniformity. Details of the search space are provided in Table~\ref{tbl:searchspace}. In particular, our formulation permits tuning an individualized rescale factor for each dimension in the RoPE embedding. To streamline the design, we opt to search over $\lambda_i$ and $\hat{n}$ as opposed to directly exploring $\mathbb{I}(\hat{\lambda}_{i},\hat{n})$, where $\hat{\lambda}_{i}=1/\lambda_i$. As detailed in Table~\ref{tbl:searchspace}, the range for $\lambda_i$ begins at 1.0 (representing simple extrapolation) and extends up to $s\times1.25$ (enabling broader interpolation than that of PI), with increments of 0.01. Here, $s$ denotes the extension ratio of the target context window.

The parameter $\hat{n}$ specifies how many of the initial token positions preserve the standard RoPE embedding, foregoing position interpolation. In practical experiments, we set $\hat{n}$ to values chosen from the set \{0, 1, 2, 4, 8, 12, 16, 20, 24, 28, 32, 64, 128, 256\} during the search. Notably, setting $\hat{n}=0$ results in every token position utilizing the optimized rescale factors.

\textbf{Evolution-based search}. The search space described in Table~\ref{tbl:searchspace} includes a vast array of positional interpolation configurations, making comprehensive search both complex and resource-intensive. For instance, when employing a $s=4\times$ extension, the total number of possible configurations escalates to $400^{128/2}\times14$ = 4$\times10^{167}$. As the extension ratio increases, the size of the search space grows exponentially. To make this process tractable, we leverage evolution search~\cite{guo2020single} and incorporate two optimization strategies that markedly enhance the efficiency of the search process. The overall methodology is summarized in Algorithm~\ref{alg:search}.

\textit{Optimized initial population generation}. Rather than randomly assigning all $P$ rescale factors for the initial population, our approach explicitly includes the RoPE rescale factors used in PI, NTK, and YaRN as distinct individuals from the outset. The other $P$-3 individuals are produced by applying random mutations to these three base rescale factors, each with probability $p$.

\textit{Monotonically non-decreasing constraint}. Once the initial population has been produced, we evaluate each individual's LLM perplexity. This involves adjusting the target LLM with the relevant RoPE rescale factors and then measuring the perplexity over the input $\mathbf{X}$. The best-performing $k$ individuals are chosen as parents for the evolutionary process. Nonetheless, because the search space is extensive, straightforward application of mutation and crossover often leads to suboptimal candidates, resulting in excessive and redundant perplexity evaluations. This issue becomes increasingly pronounced with large $L'$, as each perplexity inference incurs significant computational overhead.

To mitigate this issue, we enforce a constraint of non-decreasing monotonicity on the selected RoPE rescaling factors: $\lambda_i\le \lambda_{i+1}$. Only RoPE configurations satisfying this condition are used in the LLM perplexity analysis, which substantially limits the search space. More concretely, we stipulate that $\lambda_i$ should monotonically increase across the RoPE dimensions (with $i=0,\ldots,63$). The rationale for this monotonicity along the dimensions draws on NTK theory~\cite{jacot2018neural,tancik2020fourier,ntk}, which indicates that dimensions associated with higher frequency (i.e., lower indices) need minimal interpolation (smaller $\lambda_i$), whereas those linked to lower frequency (i.e., higher indices) can tolerate or benefit from greater interpolation (larger $\lambda_i$).

\begin{algorithm}[t]
	\small
	\caption{The search algorithm}
	\textbf{Input:} target LLM, input samples $\mathbf{X}$,   population size $P$, mutation size $N_1$, crossover size $N_2$, max iterations $\mathcal{T}$, mutate probability $p$\\

	\begin{algorithmic}[1]
		\label{alg:search}
		\STATE Initialize Top-k as an empty set;	
		\STATE Set $\text{P}_0$ by running \textit{Initialize\_population\_with\_optimization} using $P$, $\mathbf{X}$, $p$; 
		\FOR{$i$ from $1$ to $\mathcal{T}$}
		\STATE Execute \textit{Compute\_perplexity} on (LLM, $\text{P}_{i-1}$, $\mathbf{X}$);
		\STATE  Refresh Top-k with \textit{Update\_Topk} using (Top-k, $\text{P}_{i-1}$);
		\STATE Obtain $\text{P}_{mutation}$ via \textit{Mutation\_with\_mono\_constraint} applied to (Top-k, $N_1$, $p$); 
		\STATE Generate $\text{P}_{crossover}$ through \textit{Crossover\_with\_mono\_constraint} with inputs (Top-k, $N_2$);
		\STATE Form $\text{P}_i$ as the union of $\text{P}_{mutation}$, $\text{P}_{crossover}$, and Top-k;
		\ENDFOR
		\STATE Return the member in Top-k that achieves the minimum perplexity;
	\end{algorithmic}
\end{algorithm}

\textbf{8$\times$ extension without fine-tuning}. Through evolutionary search, we successfully determine non-uniform RoPE rescale factors that safeguard critical dimensions and positional encodings, thereby reducing information loss caused by interpolation. As shown in Fig.\ref{fig:non-ft}, this approach enables LLaMA2's context window to be expanded from 4k to 32k with no need for fine-tuning. Conversely, prior techniques like PI, as well as non-uniform NTK and YaRN, result in a sharp increase in perplexity after merely doubling the context size.

\subsection{Extending LLM Context Window to 2048K}
\label{sec:extendto2048k}

\textbf{Progressive extension to 2048k}. In this section, we present our approach for expanding the context window of pre-trained LLMs from the standard 4k up to beyond 2048k. Our experiments show that with non-uniform positional interpolation, it is feasible to achieve an 8$\times$ increase in context length without the need for fine-tuning. However, achieving substantially larger expansions (such as 512$\times$) necessitates fine-tuning. One possible strategy involves identifying appropriate RoPE rescale factors tailored for the target 2048k window and subsequently fine-tuning the model. Nevertheless, this process is hindered by the extremely high computational costs required. Furthermore, our observations indicate that effectively fine-tuning LLMs for such significant extension ratios proves to be challenging (refer to Appendix for details).

Luckily, LongRoPE demonstrates robust performance on both the pre-trained and the fine-tuned extended LLM. As a result, we propose a progressive and efficient strategy that can reach the target length of 2048k using only 1,000 fine-tuning iterations, with training conducted at a maximum length of 256k.

$\Diamond$ \textit{Applying LongRoPE search to expand pre-trained LLMs to 256k}. Using LLaMA2 for illustration, we perform a search targeting context window sizes of 128k and 256k. The corresponding extension ratios in this process are 32$\times$ and 64$\times$, respectively.

$\Diamond$ \textit{Fine-tuning to 256k}. Next, we adapt the pre-trained LLM for a context window of 256k. Our approach involves initially fine-tuning LLaMA2 for 400 steps employing the RoPE rescaled factors set for 128k. Once this phase is complete, we update the RoPE rescaled factors to 256k on the resulting checkpoint and continue training for another 600 steps. This staged process demonstrates greater efficiency compared to initiating fine-tuning directly at 256k.

$\Diamond$ \textit{Expanding fine-tuned extended LLM to 2048k utilizing LongRoPE search}. In the final phase, we conduct an additional search on the LLM previously fine-tuned for a 256k sequence length. This approach enables scaling the context window to a massive 2048k tokens, eliminating the need for extra fine-tuning. Consequently, the total extension reaches a factor of $512\times$.

\textbf{Recovery within the original context window}.  When expanding the context window to 2048k tokens, we observe diminished performance on tasks within the model’s initial context length. This issue has been documented in prior research on positional interpolation~\cite{pi}, which constrains positional embeddings at higher indices of the original window to occupy a much more limited subspace. Such crowding deteriorates the language model’s capabilities. Specifically, applying a 512$\times$ scale-up means that the first 4k positions are densely packed together in the new embedding space.

To address this issue, we conduct an additional round of evolution search on the extended LLM, fine-tuning the RoPE rescale factors specifically for shorter context lengths such as 4k and 8k. The maximum value permitted for the searched $\lambda$ is decreased, as shorter lengths demand less positional interpolation. While performing inference, the LLM adaptively modifies the appropriate RoPE rescale factors in real time.

 \section{Experiments}

\label{sec:exp}
\subsection{Setup}
\textbf{Evaluation Tasks and models}. We integrate LongRoPE into LLaMA2-7B and Mistral-7B architectures, assessing their efficacy from three perspectives: (1) analyzing perplexity scores for extended-context LLMs when processing lengthy texts; (2) evaluating model competency in the Passkey retrieval task, designed to assess extraction of a specified passkey amid extensive distractor content; and (3) benchmarking against standard LLM evaluation tasks using a restricted context window of 4096 tokens.

\textbf{Fine-tuning}. In the case of LLaMA2, the model is fine-tuned with a linear learning rate schedule starting at $2 \times 10^{-5}$ and a total batch size of 32 across all devices. The initial fine-tuning phase consists of 400 iterations using the Redpajama~\cite{redpajama} dataset, organized into 128k token-long sequences, each delimited by BOS and EOS markers. After completing this phase and saving a checkpoint, a further 600 training steps are performed to extend the context window up to 256k tokens. Training for the 128k sequence length is conducted on 8 A100 GPUs utilizing a distributed framework~\cite{cube}, while expanding to the 256k context window necessitates the use of 16 A100 GPUs. For Mistral, training is performed with a fixed learning rate of $1 \times 10^{-6}$ and an overall batch size of 64. Both the 128k and 256k context window models are trained according to the YaRN~\cite{yarn} methodology, comprising 400 steps using the Together Computer’s Long-Data Collections~\cite{mistraldata} at a sequence length of 16k. Training is carried out with 4 A100 GPUs.

\textbf{Search}. For target window sizes up to 256k, we adopt the following settings: $P$=64, $N_1$=$N_2$=16, $p$=0.3, and $\mathcal{T}$=40, while the top 32 candidates are retained for mutation and crossover in each search step. Perplexity evaluation employs 5 randomly chosen samples from the PG19 validation set, ensuring each sample meets or exceeds the desired context length. For windows exceeding 512k, values for population size, as well as mutation and crossover groups, are reduced by half. In this regime, perplexity is assessed using 3 random validation set samples from Pile-Books3~\cite{pile}.

\textbf{Baselines}. To achieve a context length of 2048k, we performed fine-tuning on models using context windows of 128k and 256k. As a result, we obtained LongRoPE-2048k (ft=128k) and LongRoPE-2048k (ft=256k) for LLaMA2 and Mistral, respectively. Our evaluation involves comparing these four models against leading baselines for context window extension, specifically considering open-source LLMs that have been fine-tuned following positional interpolation through PI, NTK, and YaRN methods. The baselines selected encompass Together-32k~\cite{together}, Code LLaMA~\cite{codellama}, LongLoRA-full-FT-100k~\cite{longlora}, YaRN-LLaMA, and YaRN-Mistral~\cite{yarn}.

\subsection{Main Results}

\label{sec:mainresult}
\textbf{Long sequence language modeling within 256k}. Our initial analysis focuses on benchmarking against leading extended LLMs when handling evaluation sequences up to 256k tokens. To establish the versatility of our approach, we utilize two datasets: the test splits from Proof-pile~\cite{pg19} and PG19~\cite{pile}. Perplexity is measured across a range of context lengths, employing a sliding window of size 256. For PG19, the full test set comprising 100 documents is employed. In the case of Proof-pile, adhering to the protocol of YaRN~\cite{yarn}, we randomly draw 10 samples, each with a minimum length of 128k tokens.

Table~\ref{tbl:proof-pile} and Table~\ref{tbl:pg19} present a comparison of perplexity scores for LLaMA2 and Mistral models augmented using various interpolation techniques on the Proof-pile and PG19 datasets, respectively. Two main findings emerge: \textbf{(1)} the perplexity of our extended models consistently declines as evaluation lengths increase from 4k up to 256k, demonstrating their effectiveness in utilizing longer contexts; \textbf{(2)} Despite operating within a context window that is 16$\times$ larger—a scenario that typically hinders performance on shorter contexts—our LongRoPE-2048k models surpass existing leading models within a 256k context window.

\textbf{Long sequence language modeling beyond 2000k}. To assess model performance on exceptionally lengthy texts, we utilize the Books3~\cite{pile} dataset. For efficient evaluation, we randomly sample 20 books, each longer than 2048k tokens, and apply a 256k sliding window during evaluation.

According to the results in Table~\ref{tbl:books3}, LongRoPE is capable of expanding the context window of both LLaMA2-7B and Mistral-7B models up to 2048k, while preserving perplexity levels that are similar to, or even surpass, the baseline values in the 8k-128k range. Distinct disparities emerge between the 2048k versions of LLaMA2 and Mistral: although Mistral demonstrates superior performance compared to baselines at shorter context lengths, its perplexity rises above 7 once the sequence exceeds 256k tokens. In contrast, LLaMA2 exhibits the anticipated trend of decreasing perplexity with increasing context size, with only slight upticks observed at the 1024k and 2048k marks. Notably, when using LLaMA2, training with LongRoPE-2048k at a fine-tuning length of 256k proves more effective than at 128k. This improvement stems from the reduced secondary extension ratio (8$\times$ compared to 16$\times$). Conversely, Mistral shows higher performance at a fine-tuning window size of 128k. This difference largely results from following YaRN's setup for Mistral, which applies a 16k training sequence for both 128k and 256k fine-tuning, ultimately impacting Mistral’s capacity to further scale its context window following fine-tuning.

\textbf{Passkey retrieval}. Next, we investigate the actual usable context length for generation-based tasks by employing a synthetic benchmark for passkey retrieval, originally introduced by~\cite{passkey}. In this scenario, a model must locate and output a randomly chosen five-digit number (the passkey) that is embedded within an extended passage of text. A full description of the prompt format can be found in the appendix. We conduct 10 independent runs of the passkey retrieval task, each time positioning the passkey at a randomly selected spot, drawn from a uniform distribution over the entire evaluation context range.

Fig.~\ref{fig:passkey} presents a comparative analysis of retrieval accuracy against baseline methods. The performance of existing models declines sharply, with accuracy falling to zero past the 128k token mark. Conversely, LongRoPE-LLaMA2-2048k (ft=256k) demonstrates robustness in this demanding setting, sustaining high retrieval accuracy ($\ge$90\%) across a sequence length ranging from 4k up to 2048k tokens. Similarly, LongRoPE-Mistral-2048k (ft=128k) consistently achieves perfect accuracy through 1800k, before accuracy decreases to 60\% at 2048k. This behavior is consistent with the results in Table~\ref{tbl:books3}, which indicate a modest increase in perplexity at the 2048k sequence length.

\textbf{Standard benchmarks within original context window}. We assess the performance of LongRoPE-2048k models on the initial context window by utilizing the Hugging Face Open LLM Leaderboard~\cite{huggingfaceboard} under both zero-shot and few-shot conditions. Specifically, our evaluation protocols involve the use of 25-shot ARC-Challenge~\cite{allenai:arc}, 10-shot HellaSwag~\cite{zellers2019hellaswag}, 5-shot MMLU~\cite{mmlu}, and 0-shot TruthfulQA~\cite{lin2021truthfulqa}.

As indicated in Table~\ref{tbl:commonsense}, our models perform similarly to existing systems on the original benchmark, which was intended for a limited context window, and actually surpass the original Mistral by +0.5\% on TruthfulQA. LongRoPE-LLaMA2-2048k, after being fine-tuned at 256k, experiences a minor drop in performance; nevertheless, its results for the majority of tasks are still within acceptable bounds.

\subsection{Ablation Results}

\textbf{Effectiveness of the second positional interpolation}. As part of our progressive extension method, we apply our search algorithm to perform an additional round of non-uniform positional interpolation on large language models (LLMs) that have already been fine-tuned and extended. To assess the benefits of this approach, we conduct experiments using our fine-tuned LLaMA2-256k model. This model is further extended to 512k, 1024k, and 2048k sequence lengths, utilizing both PI and YaRN for comparison. Results presented in Table~\ref{tbl:secondsearch} indicate that our non-uniform positional interpolation maintains stable perplexity across increasing sequence lengths, whereas perplexity rises sharply when applying PI and YaRN as the extension ratio grows.

\textbf{Effectiveness of recovery at shorter context lengths.} In order to address the reduction in performance at shorter context windows, we re-tuned the RoPE scaling parameters for LongRoPE-2048k using our search methodology. In this process, we set lower upper bounds for the scale factors during the search phase, thereby promoting reduced interpolation at 4k and 8k context sizes. Table~\ref{tbl:shortperformance} presents a comparison of perplexity scores for LongRoPE-LLaMA2-2048k evaluated on Proof-pile at both 4k and 8k tokens, in addition to the averaged accuracy from various LLM benchmarks. The results indicate marked gains in performance at these shorter context lengths.

\textbf{Investigation of the Two Types of Non-Uniformities}. In this section, we conduct ablation studies on the roles of both non-uniformities to determine their respective impacts on performance. Specifically, we carry out two sets of experiments: (i) We extrapolate LLaMA2-7B to shorter contexts of 16k and 32k using three approaches—PI, searching for the optimal RoPE dimension only, and searching for both non-uniformities; (ii) We further extend our model, which was fine-tuned to a length of 256k, up to 2048k, applying the same comparison strategy. Perplexity measurements are taken prior to any fine-tuning. Referring to Table~\ref{tbl:searchdimension}, adjusting the RoPE dimension non-uniformity leads to notable reductions in perplexity when contrasted with linear interpolation as implemented by PI. Adjusting non-uniformity in token position also yields perceptible gains in performance for context lengths of 16k and 32k, whereas this effect diminishes at the extreme length of 2048k. This could be due to the fact that retaining only the earliest tokens—without interpolation—loses effectiveness at such extended context lengths. This issue is left open for future investigation.

\section{Related Works}

Beyond position interpolation techniques, this section reviews alternative strategies addressed in prior research.

\textbf{Retrieval-based} strategies employ external memory components to store extensive historical context, alongside retrieval systems that access relevant documents during inference~\cite{longllama,wang2023augmenting,borgeaud2022improving}. Implementing these methods generally requires substantial architectural adjustments to LLMs. Unlike these approaches, our method involves only minimal changes to positional embeddings, offering a more efficient solution. Moreover, it supports a broader array of long-context applications, including but not limited to long document summarization and few-shot learning, extending beyond the typical scope of retrieval.

\textbf{Attention-based context window extensions}. In addition to interpolating positional embeddings, certain studies extend the input context by retaining the original LLM context window and altering the attention mechanism~\cite{han2023lminfinite, streamingllm,parallelwindow}. Central to these approaches is the introduction of new attention masks aimed at alleviating the exponential growth of attention caused by novel positions. Such approaches are compatible with, and often enhance, techniques based on positional interpolation.

\textbf{Fine-tuning based approaches} examine strategies for optimally adapting pre-trained LLMs to accommodate extended contexts by altering position embeddings. Approaches such as Code LLaMA~\cite{codellama}, LLaMA2 Long~\cite{xiong2023effective}, and ScaledRoPE~\cite{liu2023scaling} implement fine-tuning using a significantly increased RoPE base value targeted at longer sequences. In contrast, our approach provides adaptability to multiple target sequence lengths, successfully scaling to over 2M tokens. In response to the considerable GPU consumption associated with fine-tuning models for exceedingly long inputs (e.g., greater than 128k), recent techniques like LongLoRA~\cite{longlora} and PoSE~\cite{zhu2023pose} have been introduced to alleviate this computational burden. Our solution is complementary to such resource-efficient fine-tuning advancements.

\section{Conclusion}

In this paper, we introduce LongRoPE, a technique that significantly increases the context length of LLMs up to 2048k tokens, all while preserving their original performance on shorter contexts. Our approach leverages two distinct non-uniform properties within the RoPE positional embedding, utilizing an efficient evolutionary search algorithm. This results in two main advantages: it allows for optimal initialization in fine-tuning processes, and facilitates an 8$\times$ expansion of the context window even without additional fine-tuning. Furthermore, we outline a progressive extension method, where LLMs fine-tuned on sequences of 256k in length can subsequently achieve a 2048k context window without further fine-tuning efforts. Comprehensive experiments confirm the robustness and efficacy of LongRoPE. We anticipate that LongRoPE-2048k models will not only support a wide range of long-context applications, but also motivate additional advancements in this area.

\section*{Broader Impacts} This study seeks to contribute to progress within the domain of Machine Learning. While our research may have various implications for society, we do not identify any particular effects that necessitate special mention at this time.

	\balance

\end{document}